\chapter{危害国家安全罪}

% Source PDF physical page 432; see audit/source_coverage.json.
\begin{EditionOpeningPanel}
本讲很少考查，不是复习重点。
\end{EditionOpeningPanel}

\subsection*{一、间谍罪（第110条）}

1.行为方式：（1）参加间谍组织。既遂标准是加入后实施一定违法活动。（2）接受间谍组织及其代理人的任务。既遂标准是，接受任务后执行任务。成立既遂，不要求完成任务。（2025年试题）（3）为敌人指示轰击目标。既遂标准是，完成指示行为。成立既遂，不要求敌人轰击了目标。

2.罪数问题：国家机关工作人员叛逃后又参加间谍组织或者接受间谍任务的，触犯叛逃罪（第109条）和本罪，数罪并罚。

\subsection*{二、为境外窃取、刺探、收买、非法提供国家秘密、情报罪（第111条）}

本罪的“情报”须作缩小解释，只包括关系国家安全和利益、尚未公开或依照有关规定不应公开的事项。即这里的“情报”必须是关系国家安全和利益的情报，不包括一般的情报。

[总结]国家秘密类犯罪。

1.本罪与间谍罪。

（1）行为人既参加间谍组织，又为其窃取、刺探、收买、非法提供国家秘密或情报的，只定间谍罪。

（2）行为人明知对方是间谍组织，并为对方窃取、刺探、收买、非法提供国家秘密或情报的，属于“接受间谍组织的任务”，只定间谍罪。行为人如果不明知对方是间谍组织，并为对方窃取、刺探、收买、非法提供国家秘密或情报的，定本罪。

2.本罪与故意泄露国家秘密罪。

故意泄露国家秘密罪（第398条）是渎职犯罪一章的罪名。但是，国家机关工作人员和非国家机关工作人员均可以构成该罪。

两罪的区别是：二者在接受方的范围上要求不同。为境外窃取、刺探、收买、非法提供国家秘密、情报罪要求为境外机构、人员。故意泄露国家秘密罪只要求故意泄露就行，至于泄露给谁，不作要求，既可以是境内机构、人员，也可以是境外机构、人员。当故意

% Source PDF physical page 433; see audit/source_coverage.json.
泄露给境外机构、人员时，按照为境外非法提供国家秘密罪论处。\footnote{触犯为境外非法提供国家秘密罪，必然触犯故意泄露国家秘密罪，二者是法条竞合关系，前者是特殊法条，优先定前者。}

例1（2011年第5题），甲没有认识到对方是境外机构，为对方非法提供国家秘密。甲不成立为境外非法提供国家秘密罪，但是成立故意泄露国家秘密罪。

例2，通过互联网将国家秘密非法发送给境外机构的，构成为境外非法提供国家秘密罪。但是通过互联网将国家秘密只是在网上公布的，构成故意泄露国家秘密罪。

3.本罪与非法获取国家秘密罪。

非法获取国家秘密罪（第282条）是妨害社会管理犯罪一章的罪名，是指以窃取、刺探、收买方法，非法获取国家秘密的行为。

（1）非法获取国家秘密罪，主观上不能是为境外机构、组织、人员窃取、刺探、收买国家秘密，否则成立为境外窃取、刺探、收买、非法提供国家秘密、情报罪。

（2）行为人实施窃取、刺探、收买国家秘密的行为时，没有非法提供给境外机构、组织、人员的故意，但非法获取国家秘密之后，非法提供给境外机构、组织或人员的，因为侵害的法益具有同一性，仅以为境外窃取、刺探、收买、非法提供国家秘密罪论处即可，不再数罪并罚。

4.非法获取国家秘密罪与故意泄露国家秘密罪。

非法获取国家秘密后，又故意泄露该国家秘密的，重罪吸收轻罪。

\subsection*{三、叛逃罪}

第109条国家机关工作人员在履行公务期间，擅离岗位，\CJKunderwave{叛逃境外或者在境外叛逃}的，处五年以下有期徒刑、拘役、管制或者剥夺政治权利；…

掌握国家秘密的国家工作人员叛逃境外或者在境外叛逃的，依照前款的规定从重处罚。

